\section{Security Analysis}
\label{sec:security}
% ============================================================

We prove that $\Pi_{\mathsf{fpsu}}$ UC-realizes $\cF_{\mathsf{updPSU}}[f_{\cup}]$
(Fig.~\ref{fig:f-fpsu}) in the $(\cF_{\mathsf{OPRF}}, \cF_{\mathsf{OT}})$-hybrid
model against PPT semi-honest adaptive adversaries, for all four output types.

\subsection{UC Security}

\begin{theorem}[Main Security]\label{thm:main}
Assume $\SE$ is IND-CPA secure and HMAC-SHA256 is a secure PRF. Then
$\Pi_{\mathsf{fpsu}}$ UC-realizes $\cF_{\mathsf{updPSU}}[f_{\cup}]$ in the
$(\cF_{\mathsf{OPRF}}, \cF_{\mathsf{OT}})$-hybrid model against PPT semi-honest
adaptive adversaries.
\end{theorem}

\textbf{Proof structure.} The proof decomposes modularly into: (i)~per-epoch
security, inheriting from $\Pi_{\mathsf{UnionPeel}}$~\cite{PT26};
(ii)~cross-epoch forward privacy via the PRF chain (three-hybrid argument);
and (iii)~OT-masked deletion security (Lemma~\ref{lem:otmask}). The full
proof, including simulator construction, appears in
Appendix~\ref{sec:full-proofs}.

\subsection{Game-Based Forward Privacy (FP-IND)}

We supplement the UC definition with a standalone game-based notion that
directly captures the PRF-chain guarantee.

\begin{definition}[FP-IND Game]\label{def:fp-ind}
\begin{enumerate}
\item Challenger samples $b \sample \{0,1\}$, $K_0 \sample \{0,1\}^\kappa$.
\item For $t = 1$ to $T$: $K_t \leftarrow F(K_{t-1}, c)$.
\item Adversary (observing all transcripts) selects $\tau \in [T]$, receives
$K_\tau$.
\item If $b = 0$: receives real $\{K_t\}_{t < \tau}$. If $b = 1$: receives
independent $\{R_t \sample \{0,1\}^\kappa\}_{t < \tau}$.
\item Adversary outputs $b'$, wins if $b' = b$.
\end{enumerate}
$\mathsf{Adv}^{\mathsf{FP\text{-}IND}}(\kappa) = |\Pr[b' = b] - \tfrac{1}{2}|
\leq \negl(\kappa)$.
\end{definition}

\begin{corollary}[Forward Privacy]\label{cor:fp}
If $\Pi_{\mathsf{fpsu}}$ UC-realizes $\cF_{\mathsf{updPSU}}$ with adaptive
corruptions, then $\mathsf{Adv}^{\mathsf{FP\text{-}IND}}(\kappa) \leq T \cdot
\mathsf{Adv}^{\mathsf{PRF}}(\kappa)$.
\end{corollary}

\subsection{Adaptive Input Selection}

Liu et al.~\cite{LMR+26} recently showed that several published UPSI protocols
are broken by \emph{adaptive input selection}: an adversary who chooses deltas
based on prior transcripts can extract information about the honest party's
historical inputs from the protocol's stored state (e.g., OKVS encodings). We
show that F-PSU requires no modification to resist this class of attacks.

\smallskip\noindent\textbf{Why UPSI is vulnerable.}
In LMR+26's setting, the server stores an OKVS encoding of its set across
epochs. An adaptive adversary can insert elements strategically chosen to
interact with this encoding, learning membership information beyond what the
PSI output reveals. The root cause is that the server's \emph{state} depends
deterministically on its historical inputs, and adaptive queries can probe
this state.

\smallskip\noindent\textbf{Why F-PSU is immune.}
F-PSU's state has two layers of protection:
\begin{enumerate}
\item \textbf{Semantic security of encryption.} The stored IBLT is encrypted
under $\SE$ with a per-epoch key $K_t$. An adaptive adversary probing the
ciphertext learns nothing about the underlying plaintext (semantic security of
IND-CPA security of $\SE$).
\item \textbf{Forward privacy of the key chain.} The key $K_t$ is derived via a
PRF chain: $K_t \leftarrow \PRF(K_{t-1})$, with $K_{t-1}$ immediately erased.
By the FP-IND guarantee (Corollary~\ref{cor:fp}), compromising at epoch $\tau$
reveals $K_\tau$ but nothing about $\{K_t\}_{t < \tau}$. The adversary cannot
decrypt ciphertexts from prior epochs, even with adaptive oracle access.
\end{enumerate}

Consequently, the adversary's view at each epoch consists of the union output
$X_{\cup,t}$ and the size leakage $\cL_t$, exactly as specified by the ideal
functionality $\cF_{\mathsf{updPSU}}$ (Fig.~\ref{fig:f-fpsu}). Adaptive choice
of deltas changes \emph{which} elements are added or deleted, but does not
change \emph{what the protocol leaks}: the updated union and the delta sizes.
No additional game-based definition is required; the UC simulator in
Theorem~\ref{thm:main} already handles adaptive input selection by extracting
the corrupted party's effective input and forwarding it to $\cF_{\mathsf{updPSU}}$.

\smallskip\noindent\textbf{Structural contrast with PSI.}
In PSI, the output $X_{\cap,t}$ is a \emph{subset} of both inputs: an element
appears in the output only if both parties hold it. Adaptive addition of $x$
reveals $x \in Y_t$ with certainty (if $x$ appears in the new intersection).
In PSU, the output $X_{\cup,t}$ is a \emph{superset}: any element added by the
adversary appears in the union unconditionally. The adversary learns
$Y_t \setminus X_{0,t}$ (inherent to PSU), but adaptive probing yields no
additional information---the encrypted state and forward-private keys ensure
that historical data remains inaccessible regardless of input strategy.

\subsection{Security of OT-Masked Deletion}

\begin{lemma}[OT-Masked Deletion]\label{lem:otmask}
The OT-based membership check in Step~5 of $\Pi_{\mathsf{fpsu}}$ (Fig.~\ref{fig:protocol})
reveals no information beyond (i)~the per-party deletion counts $|\Delta^-_C|, |\Delta^-_S|$, and
(ii)~for each queried element $x \in \Delta^-_C$, whether $x$ belongs to $Y_{t-1}$
and $x \in \Delta^-_S$. Both are implied by the final union output $X_{\cup,t}$.
\end{lemma}

\begin{proof}[Proof sketch]
The 2-bit OT is realized via one OPRF evaluation (membership test) and one
additional OT (deletion coincidence check), both in the
$(\cF_{\mathsf{OPRF}}, \cF_{\mathsf{OT}})$-hybrid model. For Case~1 and Case~2
elements, the IBLT is unmodified, so the protocol transcript is identical to
an execution where those deletions were purely local. For Case~3 elements, the
IBLT subtraction is indistinguishable from a standard UnionPeel cell update.

A single hybrid replaces each deletion query $x$ with a random dummy element
$x' \notin X_{t-1} \cup Y_{t-1}$. The simulator can program the OT response
to match the expected output, since the final union $X_{\cup,t}$ constrains
which deletions actually removed elements. The distinguishing advantage is
bounded by $\mathsf{Adv}^{\mathsf{PRF}}(\kappa)$ per query, totaling
$|\Delta^-| \cdot \mathsf{Adv}^{\mathsf{PRF}}(\kappa) \leq \negl(\kappa)$.
\end{proof}

This lemma ensures that OT-masked deletion inherits the security of the
underlying UnionPeel protocol without introducing new attack surface. Combined
with the per-epoch UC proof (Theorem~\ref{thm:main}), the full F-PSU protocol
with OT-masked deletion UC-realizes $\cF_{\mathsf{updPSU}}[f_{\cup}]$.

% ============================================================
